% META: {"id": "1SPE-SECDEG-CO-039", "chapitre": "1SPE-SECOND-DEGRE", "type_objet": "corrige", "exercice_ref": "1SPE-SECDEG-EX-039", "capacites_codes": ["C3", "C5", "C6"], "sources_inspiration": [], "status": "generated"}
% BEGIN-VERIFY
% from sympy import *
% x = symbols('x', real=True, nonnegative=True)
% R = 15*x - x**2 / 10
% C = 200 + 5*x
% B = R - C
% B_simplified = expand(B)
% assert B_simplified == -x**2/10 + 10*x - 200
% alpha = 50
% beta = B_simplified.subs(x, alpha)
% assert beta == 50
% assert beta == 50
% # Seuil B(x) > 0 : -x^2/10 + 10x - 200 > 0 => x^2 - 100x + 2000 < 0
% B_eq = x**2 - 100*x + 2000
% Delta = 100**2 - 4*2000
% assert Delta == 2000
% r1 = (100 - sqrt(2000)) / 2
% r2 = (100 + sqrt(2000)) / 2
% assert simplify(r1 - (50 - 10*sqrt(5))) == 0
% assert simplify(r2 - (50 + 10*sqrt(5))) == 0
% import math
% assert abs(float(r1) - 27.64) < 0.1
% assert abs(float(r2) - 72.36) < 0.1
% # Q4 : B(x) = 32 => -x^2/10 + 10x - 200 = 32 => x^2 - 100x + 2320 = 0
% B_32 = x**2 - 100*x + 2320
% Delta_32 = 10000 - 4*2320
% assert Delta_32 == 720
% sols_32 = sorted(solve(B_32, x))
% assert simplify(sols_32[0] - (50 - 6*sqrt(5))) == 0
% assert simplify(sols_32[1] - (50 + 6*sqrt(5))) == 0
% assert expand(B_simplified-(-Rational(1,10)*(x-50)**2+50)) == 0
% assert B_simplified.subs(x, 28) == B_simplified.subs(x, 72) == Rational(8,5)
% # Q5 : tarif20-x/5, avec les mêmes charges200+5x.
% B2 = expand(x*(20-x/5)-C)
% assert B2 == -x**2/5+15*x-200
% alpha2 = Rational(75,2)
% beta2 = B2.subs(x, alpha2)
% assert beta2 == Rational(325,4)
% assert expand(B2-(-Rational(1,5)*(x-alpha2)**2+beta2)) == 0
% assert expand(B2-B_simplified-x*(50-x)/10) == 0
% END-VERIFY
\begin{corrige}{1SPE-SECDEG-EX-039}

\textbf{1. Recettes et bénéfice.}

\textit{a) Recette $R(x)$.}

\commentaireMarge{Recette = prix unitaire $\times$ quantité.}

$R(x) = x \cdot p(x) = x\!\left(15 - \dfrac{x}{10}\right) = 15x - \dfrac{x^2}{10}$.

\textit{b) Bénéfice $B(x) = R(x) - C(x)$.}

\[B(x) = 15x - \frac{x^2}{10} - (200 + 5x) = -\frac{x^2}{10} + 10x - 200.\] \hfill $\square$

\medskip
\textbf{2. Optimum de production.}

\textit{a) Forme canonique.}

\commentaireMarge{Compléter le carré dans $x^2-100x$.}

\[B(x)=-\frac{1}{10}(x^2-100x)-200
=-\frac{1}{10}\bigl((x-50)^2-2500\bigr)-200
=-\frac{1}{10}(x-50)^2+50.\]

\[\beta = B(50) = -\frac{2500}{10} + 500 - 200 = -250 + 500 - 200 = 50.\]

\[B(x) = -\frac{1}{10}(x - 50)^2 + 50.\]

\textit{b) Maximum.}

Pour $x = 50$ unités, le bénéfice est maximal : $B_{\max} = \mathbf{50}$~\euro.

\textit{c) Compatibilité.}

$0 \leqslant 50 \leqslant 100$ : la contrainte est satisfaite.

\medskip
\textbf{3. Seuil de rentabilité.}

\textit{a) Reformulation.}

\commentaireMarge{Multiplier par $-10$ (et inverser le sens de l'inégalité).}

$B(x) > 0 \iff -\dfrac{x^2}{10} + 10x - 200 > 0 \iff x^2 - 100x + 2000 < 0$. \hfill $\square$

\textit{b) Discriminant et racines.}

$\Delta = 100^2 - 4 \times 2000 = 10000 - 8000 = 2000$.

\[\sqrt{2000} = \sqrt{400 \times 5} = 20\sqrt{5}.\]

\[x_{1,2} = \frac{100 \pm 20\sqrt{5}}{2} = 50 \pm 10\sqrt{5}.\]

\textit{c) Intervalle de rentabilité.}

$10\sqrt{5} \approx 22{,}4$, donc $x_1 \approx 27{,}6$ et $x_2 \approx 72{,}4$.

L'entreprise est bénéficiaire exactement pour
$x\in]50-10\sqrt{5}\,;\,50+10\sqrt{5}[$.
Les bornes arrondies à l'unité sont $28$ et $72$, mais cet arrondi ne définit
pas l'intervalle exact : par exemple, $B(28)=B(72)=1{,}6>0$.

\medskip
\textbf{4. Équation de seuil $B(x) = 32$.}

\textit{a) Équation à résoudre.}

\commentaireMarge{Poser $B(x) = 32$.}

$-\dfrac{x^2}{10} + 10x - 232 = 0 \iff x^2 - 100x + 2320 = 0$.

\textit{b) Résolution.}

$\Delta = 10000 - 9280 = 720 = 144 \times 5$, donc $\sqrt{\Delta} = 12\sqrt{5}$.

\[x_{1,2} = \frac{100 \pm 12\sqrt{5}}{2} = 50 \pm 6\sqrt{5}.\]

Il y a deux solutions : $x_1 = 50 - 6\sqrt{5} \approx 36{,}6$ et $x_2 = 50 + 6\sqrt{5} \approx 63{,}4$.

\textit{Interprétation :} un bénéfice de $32$~\euro{} est atteint pour deux niveaux de production différents, l'un inférieur à l'optimal ($50$~unités) et l'autre supérieur.

\medskip
\textbf{5. Question ouverte : nouveau tarif $p'(x) = 20 - \dfrac{x}{5}$.}

\commentaireMarge{Nouveau bénéfice : $B'(x) = R'(x) - C(x)$.}

$R'(x) = x\!\left(20 - \dfrac{x}{5}\right) = 20x - \dfrac{x^2}{5}$.

$B'(x) = 20x - \dfrac{x^2}{5} - 200 - 5x = -\dfrac{x^2}{5} + 15x - 200$.

On complète le carré :
\[B'(x)=-\frac{1}{5}(x^2-75x)-200
=-\frac{1}{5}\left(\left(x-\frac{75}{2}\right)^2-\frac{5625}{4}\right)-200
=-\frac{1}{5}\left(x-\frac{75}{2}\right)^2+\frac{325}{4}.\]
Le maximum est atteint pour $x=37{,}5$, quantité admissible dans le modèle continu.

Ainsi $B'_{\max}=\frac{325}{4}=81{,}25$~\euro, contre $50$~\euro{} pour le premier tarif.

Pour comparer à une \emph{même quantité} produite et vendue, on étudie :
\[B'(x)-B(x)=-\frac{x^2}{10}+5x=\frac{x(50-x)}{10}.\]
Le tarif $p'$ donne donc un bénéfice supérieur pour $0<x<50$ ;
les deux bénéfices sont égaux en $x=0$ et $x=50$ ; le tarif $p$ est préférable
pour $50<x\leq100$. Ces conclusions comparent les deux modèles à volume fixé ;
elles ne prédisent pas comment la demande réagit à un changement de prix.

\baremeIndicatif{Q1 : 2 pts (modeliser, calculer) ; Q2 : 3 pts (calculer, raisonner) ; Q3 : 3 pts (calculer, raisonner, communiquer) ; Q4 : 2 pts (calculer, communiquer) ; Q5 : 4 pts (chercher, modeliser, communiquer)}
\end{corrige}
